% Branch B: "ordinary" outcome. Representation improves (P-1 holds) but the
% population-welfare contrast with the best static policy is not significant (P-2 fails).
\newcommand{\AbstractResults}{Under need-dependent silence, response-pooled adaptation under-served low-response residents (visit representation \vRhoSix). Population pooling with record-based screening raised their representation to \vRhoSevenR\ and their welfare by \vDWs\ (95\% CI \vDWsCI) at equal service use. Population welfare did not differ detectably from the best static policy (\vDWstat; \vDWstatCI). Record-based screening, rather than population pooling, accounted for most of the gain. Uncapped pooling amplified interpretation errors, and explicit control was more stable than an informed language-model planner.}
\newcommand{\IntroEvidence}{In a data-grounded simulation evaluated on hidden utilities, response-pooled adaptation allocated the low-response cohort a visit share below its share of need. Record-based screening and population pooling improved that cohort's representation and welfare at equal service use, without a detectable change in population welfare relative to the best static policy. Interpretation errors propagated most strongly through uncapped pooling, and the explicit controller was more stable than an informed language-model planner.}
\newcommand{\DiscussionResults}{The main effect of explicit evidence handling was distributive rather than aggregate: it changed who was served, not how much welfare the population obtained in total. Record-based screening carried most of this effect. The population-welfare comparison with the best static policy was inconclusive, so we do not claim that adaptation improves aggregate outcomes in this setting. Correction was incomplete at strong silence and absent for dissatisfaction that leaves no trace in group composition or records.}
